% Conjectura manuscript, CJ-0024 / CJ-0023: residuals of non-border-continuous
% t-norms that are (S,N)- and QL-implications.
% DRAFT FOR INDEPENDENT REVIEW. NOT REVIEWED BY A HUMAN. NOT FOR DISTRIBUTION.
% Written on 2026-11-27 (session 1 of CJ-0024) by Claude, model
% claude-opus-5-5, from research/active/CJ-0024-sn-r-nonborder-intersection/
% notes/weber-uniqueness.md (revision 2) and notes/rescaled-tnorms.md
% (revision 1). Plain LaTeX, standard packages only. Revision 1.
\documentclass[11pt]{amsart}
\usepackage{amsmath,amssymb,amsthm}
\PassOptionsToPackage{hyphens}{url}
\usepackage[colorlinks=true,allcolors=blue]{hyperref}

\newtheorem{theorem}{Theorem}
\newtheorem{lemma}[theorem]{Lemma}
\newtheorem{corollary}[theorem]{Corollary}
\newtheorem{problem}{Problem}
\renewcommand{\theproblem}{\Alph{problem}}
\theoremstyle{remark}
\newtheorem{remark}[theorem]{Remark}

\newcommand{\IWB}{I_{\mathrm{WB}}}
\newcommand{\cI}{\mathbb{I}}

\title[A second (S,N)- and QL-residual of a non-left-continuous t-norm]{Residuals of non-border-continuous t-norms that are (S,N)- and QL-implications}
\author{Conjectura project (draft; see the disclosure at the end)}
\date{27 November 2026, draft revision 1}

\begin{document}

\begin{abstract}
Baczy\'nski and Jayaram asked whether the Weber implication is the only fuzzy
implication that is simultaneously an $(S,N)$-implication and the residual of a
non-border-continuous t-norm that is not the residual of any left-continuous
t-norm, and whether it is the only QL-implication of that kind. We show that
it is not: for every left-continuous t-norm $T_0$ with a strong natural
negation and every rescaling of $[0,1]$ onto $[0,a]$ with $a$ above the fixed
point of that negation, the rescaled t-norm closed by the drastic boundary has a
residual that is both an $(S,N)$- and a QL-implication, violates the ordering
property, and differs from the Weber implication. For $a$ at or below the fixed
point the construction returns the Weber implication. Two explicit
one-parameter families, from the {\L}ukasiewicz t-norm and the nilpotent
minimum, are given with complete elementary proofs.
\end{abstract}

\maketitle

\section{The problems}

All operations are on $[0,1]$. A \emph{t-norm} (\emph{t-conorm}) is a
commutative, associative binary operation, non-decreasing in each variable,
with neutral element $1$ ($0$); a \emph{fuzzy negation} is a non-increasing
$N$ with $N(0)=1$, $N(1)=0$; a \emph{fuzzy implication} is $I$ non-increasing in
the first and non-decreasing in the second variable with $I(0,0)=I(1,1)=1$,
$I(1,0)=0$. The \emph{residual} (R-implication) of a t-norm $T$ is
$I_T(x,y)=\sup\{t\in[0,1] : T(x,t)\le y\}$, with no continuity assumed; an
\emph{$(S,N)$-implication} is $I_{S,N}(x,y)=S(N(x),y)$; a \emph{QL-implication}
is a fuzzy implication of the form $S(N(x),T(x,y))$. $T$ is
\emph{border-continuous} if it is continuous at every point of
$[0,1]^2\setminus\,]0,1[^2$. The \emph{ordering property} (OP) is
$I(x,y)=1\iff x\le y$. The Weber implication is $\IWB(x,y)=1$ for $x<1$,
$\IWB(1,y)=y$; it is the residual of the drastic product $T_D$ and equals
$I_{S,N_{D2}}$ for every t-conorm $S$. We write $\cI_T$, $\cI_{T_{LC}}$,
$\cI_{S,N}$, $\cI_{QL}$ for the R-implications of all t-norms, of
left-continuous t-norms, the $(S,N)$- and the QL-implications.

The following questions are quoted from \cite{BJ2013} (Section~4, p.~202),
where they are said to remain open; the second is also Problem~9.4(ii) of
\cite{BJ2010}, and the first is Problem~4.8.1 of \cite{BJ2008book}, quoted in
\cite[\S3.1.3]{BJ2013} as ``is the following equality true
$\cI_{S,N}\cap(\cI_T\setminus\cI_{T^*})=\{\IWB\}$?''.

\begin{problem}[{\cite[Problem 4.2]{BJ2013}}]\label{pr:sn}
(i) Is there a fuzzy implication $I$, other than the Weber implication
$\IWB$, which is both an $(S,N)$-implication and an R-implication which is
obtained from a non-border continuous t-norm and cannot be obtained as the
residual of any other left-continuous t-norm? (ii) If the answer to the above
question is in the affirmative, characterize the above non-empty
intersection.
\end{problem}

\begin{problem}[{\cite[Problem 4.5(ii)]{BJ2013}}]\label{pr:ql}
Is the Weber implication $\IWB$ the only QL-implication that is also an
R-implication obtained from a non-left continuous t-norm? If not, give other
examples from the above intersection and hence, characterize the non-empty
intersection $\cI_{QL}\cap\cI_T$.
\end{problem}

\begin{remark}[Reading of Problem~\ref{pr:ql}]
Read literally, the first question of Problem~\ref{pr:ql} is answered in the
sources themselves: the non-left-continuous nilpotent minimum ($0$ for
$x+y<1$, $\min$ otherwise) has the Fodor implication as residual
\cite[p.~189]{BJ2013}, and the Fodor implication is a QL-implication
\cite[Table~6, Proposition~6.4]{BJ2010}. We therefore read it, as the
monograph form of Problem~\ref{pr:sn} does, as asking for QL-implications in
$\cI_T\setminus\cI_{T_{LC}}$. The results below are claimed only under this
reading.
\end{remark}

\section{Two lemmas}

The following construction is classical (\cite[Corollary~1.8]{KMP}, as used in
\cite[\S4, (5)--(6)]{JBM2011}); proofs are included for completeness.

\begin{lemma}\label{lem:drastic}
Let $M:[0,1]^2\to[0,1]$ be commutative, associative, non-decreasing, with
$M\le\min$. Then $T(x,y)=\min(x,y)$ if $\max(x,y)=1$ and $T(x,y)=M(x,y)$
otherwise is a t-norm, and for $x<1$
\[
  I_T(x,y)=\sup\bigl(\{t\in[0,1[\;: M(x,t)\le y\}\cup\{1: x\le y\}\bigr),
  \qquad I_T(1,y)=y .
\]
\end{lemma}

\begin{proof}
Monotonicity: for $x<1$, $T(x,t)=M(x,t)\le\min(x,t)\le x=T(x,1)$. For
associativity, if one of $x,y,z$ equals $1$ both sides reduce to $T$ of the
other two; if all are $<1$ then $T(x,y)=M(x,y)<1$, $T(y,z)<1$, and both sides
equal $M(M(x,y),z)=M(x,M(y,z))$. The residual formula is the definition, with
$T(x,1)=x$.
\end{proof}

\begin{lemma}\label{lem:lc}
If $T$ is a t-norm with $T(x,\cdot)$ left-continuous at $1$ and
$I_T(x,y)=1$, then $x\le y$. Hence an implication $I$ with $I(x_0,y_0)=1$ for
some $x_0>y_0$ is not the residual of any left-continuous t-norm.
\end{lemma}

\begin{proof}
If $T(x,1)\le y$ we are done. Otherwise $\sup\{t:T(x,t)\le y\}=1$ is not
attained, so by monotonicity $T(x,s)\le y$ for all $s<1$, and
$x=T(x,1)=\lim_{s\to1^-}T(x,s)\le y$.
\end{proof}

\section{The general construction}

\begin{theorem}\label{thm:C}
Let $T_0$ be a left-continuous t-norm whose natural negation
$N_0(v)=\sup\{t: T_0(v,t)=0\}$ is strong, $e$ its fixed point, and $S_0$ the
$N_0$-dual t-conorm. Let $0<a<1$ and let $\psi:[0,1]\to[0,a]$ be an increasing
bijection. Put $T(x,y)=\min(x,y)$ if $\max(x,y)=1$,
$T(x,y)=\psi^{-1}(T_0(\psi x,\psi y))$ otherwise;
$N(x)=\psi^{-1}(\min(a,N_0(\psi x)))$ for $x<1$, $N(1)=0$;
$S(p,q)=\psi^{-1}(\min(a,S_0(\psi p,\psi q)))$. Then:
\begin{enumerate}
\item $T$ is a t-norm, $S$ a t-conorm, $N$ a fuzzy negation;
\item $I_T(x,y)=\psi^{-1}(\min(a,I_{T_0}(\psi x,\psi y)))$ for $x<1$;
\item $I_T(x,y)=S(N(x),y)=S(N(x),\min(x,y))$ for all $x,y$;
\item if $a\le e$, then $I_T=\IWB$;
\item if $e<a<1$, then $I_T(x,0)=1$ for $0<x\le\psi^{-1}(N_0(a))$, so $I_T$
  is not the residual of any left-continuous t-norm; $T$ is not
  border-continuous; and $I_T\ne\IWB$.
\end{enumerate}
\end{theorem}

\begin{proof}
We use: (F1) $T_0(u,t)\le v\iff t\le I_{T_0}(u,v)$ (left-continuity);
(F2) $I_{T_0}(u,v)=S_0(N_0(u),v)$ \cite[Theorem~5.13]{BJ2010}; (F3)
$I_{T_0}(u,u)=1$, $I_{T_0}(u,0)=N_0(u)$; (F4) $N_0(u)\ge a\iff u\le N_0(a)$,
and $a>e\iff N_0(a)<e<a$; (F5) $T_0(u,v)=0$ when $u\le N_0(v)$, since
$T_0(v,N_0(v))=0$ by (F1).

(1) $M=\psi^{-1}\circ T_0\circ(\psi\times\psi)$ is well defined, being
conjugate to the restriction of $T_0$ to $[0,a]^2$, which $T_0$ maps into
$[0,a]$; it is associative, commutative, non-decreasing, $M\le\min$; apply
Lemma~\ref{lem:drastic}. With $U=\min(a,S_0)$ on $[0,a]^2$ one has
$U(U(u,v),w)=\min(a,S_0(S_0(u,v),w))$ (both sides equal $a$ if
$S_0(u,v)\ge a$), a symmetric expression; so $U$ is associative, and $S$,
conjugate to $U$, is a t-conorm. $N$ is clearly a fuzzy negation.

(2) For $x<1$ and $t<1$: $M(x,t)\le y\iff \psi t\le r:=I_{T_0}(\psi x,\psi y)$
by (F1). If $r\ge a$ the set is $[0,1[$, supremum $1=\psi^{-1}(a)$; otherwise
it is $[0,\psi^{-1}(r)]$. If $x\le y$ then $r=1$. Lemma~\ref{lem:drastic}
gives the formula.

(3) For $x=1$ both sides equal $y$. For $x<1$, with $u=\psi x$, $v=\psi y$:
$S(N(x),y)=\psi^{-1}(\min(a,S_0(\min(a,N_0u),v)))=
\psi^{-1}(\min(a,S_0(N_0u,v)))$ (both are $1$ if $N_0u\ge a$), which is
$I_T(x,y)$ by (F2) and (2). If $y\ge x$, $S(N(x),x)=\psi^{-1}(\min(a,I_{T_0}(u,u)))=1$ by
(F3); if $y<x$, $\min(x,y)=y$.

(4) If $a\le e$, then for $x<1$, $N_0(\psi x)\ge N_0(a)\ge e\ge a$, so
$I_T(x,0)=1$ and $I_T(x,\cdot)\equiv1$.

(5) For $0<x\le\psi^{-1}(N_0(a))$, $N_0(\psi x)\ge a$, so $I_T(x,0)=1$;
apply Lemma~\ref{lem:lc}. For the same $x$ and $t<1$,
$T_0(\psi x,\psi t)\le T_0(\psi x,a)=0$ by (F5), so $T(x,t)=0$ while
$T(x,1)=x>0$. For $N_0(a)<\psi x<a$, $I_T(x,0)=\psi^{-1}(N_0(\psi x))<1$.
\end{proof}

\section{Two explicit families}

Let $1<c<2$ and $N_c(x)=\min(1,c-x)$ for $x<1$, $N_c(1)=0$. Taking
$\psi(x)=x/c$ (so $a=1/c>1/2=e$) in Theorem~\ref{thm:C}:

\begin{corollary}\label{cor:A}
$T_c(x,y)=\max(0,x+y-c)$ for $x,y<1$, $\min(x,y)$ otherwise, is a
non-border-continuous t-norm with residual $I_c(x,y)=\min(1,c-x+y)$ ($x<1$),
$I_c(1,y)=y$, and $I_c(x,y)=\min(1,N_c(x)+y)=\min(1,N_c(x)+\min(x,y))$.
($T_0=T_{\mathbf L}$.)
\end{corollary}

\begin{corollary}\label{cor:B}
$T'_c(x,y)=\min(x,y)$ if $\max(x,y)=1$ or $x+y>c$, and $0$ otherwise, is a
non-border-continuous t-norm with residual $J_c(x,y)=1$ ($x\le y$),
$J_c(x,y)=\max(\min(1,c-x),y)$ ($y<x<1$), $J_c(1,y)=y$, and
$J_c(x,y)=S^*_c(N_c(x),y)=S^*_c(N_c(x),\min(x,y))$ with the t-conorm
$S^*_c(a,b)=\max(a,b)$ for $a+b<c$, $1$ otherwise. ($T_0=T_{nM}$.)
\end{corollary}

Both corollaries are also proved directly, without Theorem~\ref{thm:C}, in the
project note cited below, and both were checked in exact rational arithmetic on
grids (residuals bracketed by exact bisection). In each family $I_c(c-1,0)=1$,
$I_c\ne\IWB$, and different $c$ give different implications; $I_c\ne J_c$
although both have the natural negation $N_c$.

\begin{corollary}
The answer to Problem~\ref{pr:sn}(i) is yes, and the equality
$\cI_{S,N}\cap(\cI_T\setminus\cI_{T_{LC}})=\{\IWB\}$ is false: the
intersection contains continuum many elements. Under the reading of the
remark above, the answer to the first question of Problem~\ref{pr:ql} is no.
\end{corollary}

\begin{remark}
Theorem~6.5 of \cite{BJ2008} (Theorem~5.14 of \cite{BJ2010}) is printed for a
right-continuous t-conorm $S$ and an arbitrary t-norm $T$. Corollary~\ref{cor:B}
($S=S^*_c$ is right-continuous with $N_{S^*_c}=N_c$, not strict) contradicts
it in that generality, as does $\IWB=I_{S_M,N_{D2}}=I_{T_D}$, noted in
\cite[Remark~6.4(i)]{BJ2008} immediately before it; the remark stresses that
left-continuity of $T$ is essential, which suggests that this hypothesis was
intended.
\end{remark}

\section*{Status, evidence and disclosure}

The characterisation parts of both problems remain open. A literature audit
(27~November 2026; the sources listed below, the citation lists of
\cite{BJ2008,BJ2010} and of the 2015 handbook chapter by Baczy\'nski,
Jayaram, Massanet and Torrens, screened by title) found no earlier answer and
no occurrence of these families; the 2015 chapter and \cite{BJ2008book} were
not accessible, so novelty is not established. The project files are
\texttt{research/active/CJ-0024-sn-r-nonborder-intersection/} (notes, reviews)
and \texttt{experiments/CJ-0024-sn-r-nonborder-intersection/} (exact checks).
This draft, the construction and the proofs were produced by an AI model
(Claude, \texttt{claude-opus-5-5}) in the Conjectura project directed by
Gustavo Bernardo Brito; the reviews so far are AI reviews; no human has
checked the mathematics.

\begin{thebibliography}{9}
\bibitem{BJ2008} M.~Baczy\'nski, B.~Jayaram, $(S,N)$- and R-implications: a
state-of-the-art survey, Fuzzy Sets Syst.\ 159 (2008) 1836--1859.
\bibitem{BJ2008book} M.~Baczy\'nski, B.~Jayaram, Fuzzy Implications,
Studies in Fuzziness and Soft Computing 231, Springer, 2008.
\bibitem{BJ2010} M.~Baczy\'nski, B.~Jayaram, QL-implications: some
properties and intersections, Fuzzy Sets Syst.\ 161 (2010) 158--188.
\bibitem{BJ2013} M.~Baczy\'nski, B.~Jayaram, Fuzzy implications: some
recently solved problems, in: Advances in Fuzzy Implication Functions, Studies
in Fuzziness and Soft Computing 300, Springer, 2013, 177--204.
\bibitem{JBM2011} B.~Jayaram, M.~Baczy\'nski, R.~Mesiar, R-implications and
the exchange principle: a complete characterization, in: Proc.\ EUSFLAT 2011,
Atlantis Press, 223--229.
\bibitem{KMP} E.~P.~Klement, R.~Mesiar, E.~Pap, Triangular Norms, Kluwer,
2000 (cited after \cite{JBM2011}; not consulted).
\end{thebibliography}

\end{document}
